Honest caveats, aggregated from the implementation chapters so a
reviewer can see them in one place. Each threat below lists its source
(how we noticed it), its likely impact on the reported results, and the
mitigation we did or did not apply. The structure follows the MDPI
methodology-article convention of explicit construct / internal /
external / conclusion validity threats, adapted to an embedded-systems
measurement context.

\begin{center}\rule{0.5\linewidth}{0.5pt}\end{center}

\section{1. Construct validity --- are we measuring what we
claim?}\label{construct-validity-are-we-measuring-what-we-claim}

\subsection{1.1 ``FPS'' is always a derived statistic, not an observed
one}\label{fps-is-always-a-derived-statistic-not-an-observed-one}

Every reported frame-per-second number in this paper is computed as
\texttt{1000\ /\ mean\_total\_ms} from a raw CSV of per-iteration
wall-clock measurements. We never sample frame counts over a long
window. The two forms are mathematically equivalent only if the
distribution of iteration times is well-behaved (low skew, bounded
tail). In practice the distribution is slightly right-skewed at both 15
fps and 30 fps sensor rates (typical \texttt{σ} of 1.5--3 ms on a
60--150 ms mean); the derived FPS is therefore within 0.1 \% of a
hypothetical window-count measurement. See \url{statistical_notes.md}
§``FPS derivation''.

\textbf{Mitigation}: report both \texttt{total\_ms} mean and derived FPS
in every result table; a reviewer can re-derive the other.

\subsection{1.2 Per-stage timings add Python dispatch overhead that
varies}\label{per-stage-timings-add-python-dispatch-overhead-that-varies}

\texttt{select\_ms}, \texttt{to\_tensor\_ms}, \texttt{detect\_ms} are
measured with \texttt{ticks\_ms()} snapshots around the MicroPython
dispatch. Each snapshot adds \textasciitilde100 µs of Python interpreter
overhead to the measured stage. At the 1 ms tick resolution of the
counter this is below the noise floor and does not affect the reported
numbers, but it is worth flagging for a reviewer who adds their own
instrumentation: we never claim sub-millisecond precision on any
host-side stage timer.

The one stage whose duration is reported from the firmware itself
(\texttt{invoke\_ms} --- returned by \texttt{sentai.tpu.invoke()}) has
no Python overhead in its measurement. That is why \texttt{invoke\_ms}
has the tightest reproducibility in the cross-session table (Table 4 of
\url{evaluation.md}).

\subsection{1.3 The alternating sweep measures a worst-case scheduling
policy}\label{the-alternating-sweep-measures-a-worst-case-scheduling-policy}

Experiment E16 / E18 sweep C flips the MUX on \emph{every} iteration.
Nothing in the shipping firmware requires the application layer to
alternate this aggressively. The 6.86 FPS ``alternating'' number is
therefore a \textbf{lower bound} on application throughput, not an
expected operating point. Users adopting the
\texttt{sentai.camera.ratio(a,\ b)} scheduler for asymmetric capture
will see effective FPS interpolate between the fixed (\textasciitilde16
FPS) and the alternating (\textasciitilde6.9 FPS) extremes as a function
of the switch density.

\begin{center}\rule{0.5\linewidth}{0.5pt}\end{center}

\section{2. Internal validity --- could something else explain the
results?}\label{internal-validity-could-something-else-explain-the-results}

\subsection{2.1 Single-board
measurements}\label{single-board-measurements}

Every result in this paper was measured on one physical SentAI v1.0
board. We did not swap OV5640 camera modules, did not change the host
USB cable or port between sessions, and did not run on a different
instance of the same hardware revision.

Consequences: - A between-sample-board deviation (e.g.~sensor silicon
variation causing different AEC/AGC convergence) would not be visible in
our measurements. - The 2 ms residual cam1 − cam0 asymmetry reported
post-Fix B may or may not reproduce on a different unit --- we cannot
tell.

\textbf{Mitigation}: the measurement framework is fully open and
reproducible (see \url{artifact.md}); a follow-up study on a second
board would directly test this.

\subsection{2.2 Static-scene confound}\label{static-scene-confound}

All measurements were captured against a static scene documented by the
before/after JPEGs in every session folder (see
\url{experimental_setup.md} §5). The scene happens to contain zero
instances of the model's target class, which means: -
\texttt{num\_detections} is 0 on every frame, so NMS takes a consistent
(and fast) path. - We do not observe the tail behaviour of the
post-processing pipeline when detection counts spike.

\textbf{Mitigation}: the per-iteration CSVs record
\texttt{num\_detections} so a future measurement against a busy scene
can compare directly. The scene JPEGs are checked into
\href{../experiments/}{\texttt{../experiments/}} so a reviewer can
confirm the scene visually.

\subsection{2.3 Thermal state across
sessions}\label{thermal-state-across-sessions}

Each of our sessions takes 3--10 seconds of wall time. The RT1176 does
not thermal-throttle in that window, and we did not observe any drift in
the per-iteration times within a session. But the paper's cross-session
comparison (Table 4 of \url{evaluation.md}) compares sessions that may
have been taken with different ``warm-up'' exposure of the host
environment --- e.g.~the host USB subsystem's buffering state can
differ.

The ≤ 1 ms cross-session agreement observed in Table 4 is evidence
against any material thermal or host-state confound at the sub- second
measurement granularity we use. We do not claim robustness for
multi-minute sustained-stress measurements.

\subsection{2.4 Firmware build drift across a reported
comparison}\label{firmware-build-drift-across-a-reported-comparison}

The cross-session comparison in Table 4 of \url{evaluation.md} spans
three firmware builds (pre- refactor, post-refactor confirmation,
post-NASA-JPL-review). Each build was flashed with
\texttt{scripts/flashtool.py\ -e\ sentai\_runtime} (no \texttt{-\/-ram};
true persistent flash) and rebooted cleanly before the measurement. The
\texttt{build\_version.h} number in each session's \texttt{summary.txt}
documents which build produced the data.

The intent of the comparison is \emph{verifying performance neutrality}
of the review refactor, not isolating a specific change. We do not claim
that any single NASA-JPL fix in isolation is responsible for the 0.3 ms
shift in cam0 total between \texttt{s043} and \texttt{s045}; within
noise, it is not.

\subsection{\texorpdfstring{2.5 The \texttt{drain=1} knob is retained
but not
recommended}{2.5 The drain=1 knob is retained but not recommended}}\label{the-drain1-knob-is-retained-but-not-recommended}

The shipping default is \texttt{switch\_drain(2)}. Moving to
\texttt{switch\_drain(1)} was empirically (see \url{cam_switch.md}
§``Known limitations''):

\begin{itemize}
\tightlist
\item
  \textbf{Not a timing win}: the \texttt{cam-\textgreater{}GetRawFrame}
  call after the drain loop blocks long enough to absorb the one-frame
  saving. Observed: 201 ms per iteration at both \texttt{drain=1} and
  \texttt{drain=2} in
  \href{../experiments/s041_e17_eof_check/}{\texttt{s041\_e17\_eof\_check}}.
\item
  \textbf{Not visually clean}: residual sensor-side artefacts (AEC/AGC
  convergence on the newly selected OV5640) remain even after the
  mid-buffer seam is eliminated by the flip-on-EOF fix. Not every frame
  is affected; a thumbnail-level review misses it.
\end{itemize}

Retaining the knob as an experimentation hook has a small cost (the
\texttt{switch\_drain\_set/\_get} accessors in the MicroPython module);
the value of being able to reproduce the limitation is greater than the
cost of carrying the code. No production configuration should use
\texttt{switch\_drain(1)} on this hardware.

\begin{center}\rule{0.5\linewidth}{0.5pt}\end{center}

\section{3. External validity --- where do the results
generalise?}\label{external-validity-where-do-the-results-generalise}

\subsection{3.1 Model-dependent numbers}\label{model-dependent-numbers}

Every reported invoke time (30 ± 3 ms in E18) and every per-stage total
is specific to the YOLOv5-enhanced single-class model documented in
\url{experimental_setup.md} §4. A heavier backbone, a larger head, or
more classes would change \texttt{invoke\_ms} without changing the
per-switch tax (which is camera- topology bound). A lighter model would
do the opposite --- making the switch tax a larger fraction of the
total.

When a reader adapts these results to a new model, the rule is: - Switch
tax per flip (\textasciitilde83 ms) is \textbf{stable} against model
choice. - TPU invoke time scales \textbf{\textasciitilde linearly} with
input-tensor pixel count and roughly with FLOPs on EdgeTPU. - PXP resize
+ quant scales with the output size (fixed at 512×512 in our
measurements) and is \textbf{not} a strong function of sensor native
resolution.

See \url{cam_switch.md} §``What E18 tells us about application design''
for the predictive arithmetic linking these knobs to effective FPS.

\subsection{3.2 Sensor-topology-dependent
numbers}\label{sensor-topology-dependent-numbers}

The 83 ms per-switch tax is specific to the single-MIPI-lane +
analogue-MUX topology of the SentAI board. Platforms with two
independent MIPI-CSI2 receivers (or an MIPI-bridge chip with virtual
channel support --- see \url{related_embedded_inference.md}) would not
pay this tax at all. The \emph{method} --- flip-on-EOF timing, ratio
scheduler, fault counters --- generalises; the \emph{numbers} do not.

\subsection{3.3 Resolution-dependent
expectations}\label{resolution-dependent-expectations}

All \texttt{evaluation.md} numbers are at 512×512 logical resolution.
Experiments also ran at VGA (640×480) and QVGA (320×240) and are
documented inline in \url{cam_switch.md} §``Per-camera resolution''.
Summary: - Fixed-camera FPS is \textbf{sensor-rate-bound} above
\textasciitilde16 FPS; changing resolution does not raise the ceiling
meaningfully. - Alternating FPS improves noticeably at QVGA
(\textasciitilde10 FPS) because PXP + JPEG costs scale with pixel count.
- VGA is \textbf{slightly worse} than 512×512 for JPEG encoding because
it has more pixels (307 k vs 262 k).

\subsection{3.4 Hardware revision scope}\label{hardware-revision-scope}

Every result applies to \textbf{SentAI board v1.0 only}. A future board
revision that rewires the MUX, adds a second CSI receiver, or changes
the sensor modules would invalidate the constants in this paper but
should preserve the \emph{framework} (session-based measurement,
flip-on-EOF as a safe MUX-transition primitive, fault counters as a
measurement-integrity check).

\begin{center}\rule{0.5\linewidth}{0.5pt}\end{center}

\section{4. Conclusion validity --- are the statistics
defensible?}\label{conclusion-validity-are-the-statistics-defensible}

\subsection{4.1 Sample size}\label{sample-size}

\texttt{n\ =\ 20} (E15) and \texttt{n\ =\ 40} (E16, E18) with the first
sample dropped. With \texttt{σ\ ≈\ 2\ ms} on a \textasciitilde60 ms mean
this gives a 95 \% confidence interval of \texttt{±0.9\ ms}, well below
the effect sizes we report (+9.5 ms from eDMA, +83 ms switch tax, −64 ms
asymmetry collapse). See \url{statistical_notes.md} for the derivation.

\subsection{4.2 No formal hypothesis
testing}\label{no-formal-hypothesis-testing}

We do not compute p-values. All of our ``effects'' are at least an order
of magnitude larger than the measurement noise; a t-test against a null
of ``no change'' would produce \texttt{p\ \textless{}\ 10⁻¹⁰} and not
add information. Where we report a null result (Fix A having no
measurable effect), we say so explicitly and show the raw numbers so a
reviewer can confirm.

\subsection{4.3 Cross-session comparisons rely on held-constant
covariates}\label{cross-session-comparisons-rely-on-held-constant-covariates}

Every cross-session table in this paper holds model, scene, resolution
and sensor configuration constant, varying only the firmware build or an
A/B runtime flag. If any of those covariates had drifted, it would show
up as a shift in the fixed-camera baseline --- which empirically
remained at 62.6 ± 0.4 ms across
\href{../experiments/s043_e18_headtail_drain2/}{\texttt{s043}},
\href{../experiments/s044_e18_headtail_drain2/}{\texttt{s044}}, and
\href{../experiments/s045_e18_post_refactor/}{\texttt{s045}} (see Table
4 of \url{evaluation.md}). That agreement is the strongest available
evidence that the comparisons isolate the variable under test.

\subsection{4.4 Single-run sessions excluded from
statistics}\label{single-run-sessions-excluded-from-statistics}

Sessions \texttt{s016}--\texttt{s020} and \texttt{s027}--\texttt{s030}
contain a single data point each (n = 1) and are not used for any mean
comparison in this paper. They are retained in the archive (Appendix C
of \href{../experiments/README.md}{experiments/README.md}) because they
document the progression of single-shot pre/post-eDMA checks before the
x20 repetition protocol stabilised; they do not contribute statistical
weight to the headline results.

\begin{center}\rule{0.5\linewidth}{0.5pt}\end{center}

\section{5. What we did not do, and why it is not a
threat}\label{what-we-did-not-do-and-why-it-is-not-a-threat}

For transparency, the following were considered and explicitly declined
in this iteration:

{\def\LTcaptype{none} % do not increment counter
\begin{longtable}[]{@{}
  >{\raggedright\arraybackslash}p{(\linewidth - 4\tabcolsep) * \real{0.3333}}
  >{\raggedright\arraybackslash}p{(\linewidth - 4\tabcolsep) * \real{0.3333}}
  >{\raggedright\arraybackslash}p{(\linewidth - 4\tabcolsep) * \real{0.3333}}@{}}
\toprule\noalign{}
\begin{minipage}[b]{\linewidth}\raggedright
Not done
\end{minipage} & \begin{minipage}[b]{\linewidth}\raggedright
Why not
\end{minipage} & \begin{minipage}[b]{\linewidth}\raggedright
When it would matter
\end{minipage} \\
\midrule\noalign{}
\endhead
\bottomrule\noalign{}
\endlastfoot
Cross-device measurement & single prototype hardware & production
deployment across N units \\
Thermal stress over minutes & 3-s measurement windows don't throttle &
always-on field deployment \\
Power / current measurement & no instrumentation on the test rig &
battery-bounded mission profile \\
FSIN master/slave sensor sync & hardware rework required, ≤ 2 ms
residual asymmetry not worth the silicon change & if directional
asymmetry became the dominant error source \\
60 fps 720p mode & NXP driver binning-mode init missing; CSI-2 did not
lock on our probe & a use case forced by a model whose per-frame cost
fits a 17 ms budget \\
Virtual-channel-based parallel capture & requires a MIPI-bridge chip not
on the board & a dual-sensor application that cannot tolerate any switch
tax \\
\end{longtable}
}

Each of these is a future-work item (\url{future_work.md}) rather than a
concealed threat to the results as reported.
